

	

\thispagestyle{empty}~~\vspace{-5cm}
%\begin{changemargin}{-1cm}{-1cm}

\BgThispage

%\thispagestyle{empty}~~\vspace{-0.2cm}

%	\newpage


%%%%%%%%%%%%%%%%%%%%%%%%%%%%%%%%%%%%%%%%%%%%%%%%%%%%%%%%%%%%%%%%%%%%%%%%%%%%%%%%%%%%%%%%%%%%%% 
%%%%%%                        Va Una Pagina                                                    %%
%%%%%%%%%%%%%%%%%%%%%%%%%%%%%%%%%%%%%%%%%%%%%%%%%%%%%%%%%%%%%%%%%%%%%%%%%%%%%%%%%%%%%%%%%%%%%%


%\begin{center}
%	
\includegraphics[width=0.2\linewidth, height=0.2\textheight]{Upto}
%\LAGE{ULA}
%\end{center}


\vspace{2.0cm}{
{\hspace{-0.3cm}\rotatebox[origin=c]{13}{ {{\bf\LARGE UNIVERSIDAD DE LOS ANDES}}}
 }}

%\hspace{-3.75cm}
\vspace{-1.25cm}{
\hspace{0.15cm}\rotatebox[origin=c]{13}{{\bf \large{MÉRIDA, VENEZUELA} }}} 
%\rotatebox[origin=c]{33}{texto siguiente esta rotado 33 grados respecto al centro}
%\vspace{2cm}
\vspace{13cm} %  17.3cm

%\end{flushright}

%\begin{flushleft}



\vspace{-2.75cm}{}{\hspace{3.60cm}\vspace{-5cm}
\includegraphics[scale=0.25,angle=17]{Celciec}}   %  \vspace{0.25cm}


{\vspace{3.5cm}{}}

{\hspace{4.20cm}\rotatebox[origin=c]{17}{\bf \Large{Autor: Linis B. Guerrero P}}}

%{\vspace{-3.00cm}\hspace{4.5cm}\rotatebox[origin=c]{17}{\bf \Large{Traducido: Prof. Linis B. Guerrero P.}}} 
%\end{flushleft}



	
%\vspace{2.0cm}{
%	{\hspace{-0.3cm}\rotatebox[origin=c]{13}{ {{\bf\LARGE UNIVERSIDAD DE LOS ANDES}}}}}


\thispagestyle{empty}~~\vspace{-0.2cm}


\newpage


\backgroundsetup{contents=Borrador, opacity=0.45, scale=10, color=blue, angle=45}



%%%%%%%%%%%%%%%%%%%%%%%%%%%%%%%%%%%%%%%%%%%%%%%%%%%%%%%%%%%%%%%%%%%%%%%%%%%%%%%%%%%%%%%%%%%%%% 
%%%%%%                        Va Una Pagina                                                    %%
%%%%%%%%%%%%%%%%%%%%%%%%%%%%%%%%%%%%%%%%%%%%%%%%%%%%%%%%%%%%%%%%%%%%%%%%%%%%%%%%%%%%%%%%%%%%%%

%	\hspace{2cm}

\begin{center}
	
\includegraphics[width=0.2\linewidth, height=0.2\textheight]{ULA}
\end{center}


\begin{center}
	{\ {\bf \vspace{0.1cm}  UNIVERSIDAD DE LOS ANDES\\	MÉRIDA\\
			\vspace{0.1cm} 	 VENEZUELA\\
	}}
\end{center}



\vskip 0.5cm \rule{\linewidth}{2.3pt}\\%[1cm]
\begin{center}
	{\  \LARGE  {\bf \vspace{0.2cm} Un curso intensivo de Python 3 \\			 
			\vspace{0.2cm} Una introducción rápida a la programación con Python 3	}} \\ 		
\end{center}

\vskip 0.5cm \rule{\linewidth}{2.3pt}\\[0.5cm]  

{\bf {\vspace{-1cm} Una Guía rápida a la programación con Python 3}}




\vspace{2cm} 
\begin{center}
	{\ {\bf \vspace{0.2cm}  	Autor: Linis B. Guerrero P.\\
			Traducido: Prof. Linis B. Guerrero P.\\
			\vspace{0.2cm} 	
	}}
\end{center}


\begin{center}
	{\ {\bf \vspace{0.2cm}  	Mérida, Venezuela \\
			Septiembre, 2023 \\
			\vspace{0.2cm} 
	}}
\end{center}






\thispagestyle{empty}~~\vspace{-0.2cm}



%%%%%%%%%%%%%%%%%%%%%%%%%%%%%%%%%%%%%%%%%%%%%%%%%%%%%%%%%%%%%%%%%%%%%%%%%%%%%%%%%%%%%%%%%%%%%% 
%%%%%%                        Va Una Pagina                                                    %%
%%%%%%%%%%%%%%%%%%%%%%%%%%%%%%%%%%%%%%%%%%%%%%%%%%%%%%%%%%%%%%%%%%%%%%%%%%%%%%%%%%%%%%%%%%%%%%








\begin{center}
	
\includegraphics[width=0.2\linewidth, height=0.2\textheight]{ULA}
\end{center}


\begin{center}
	{\ {\bf \vspace{0.1cm}  UNIVERSIDAD DE LOS ANDES\\	MÉRIDA\\
			\vspace{0.1cm} 	 VENEZUELA\\
	}}
\end{center}


\vskip 0.5cm \rule{\linewidth}{3pt}\\%[1cm]
\begin{center}
	{\  \LARGE  {\bf \vspace{0.2cm} Un curso intensivo de Python 3 \\			 
			\vspace{0.2cm} Una introducción rápida a la programación con Python 3	}} \\ 		
\end{center}

\vskip 0.5cm \rule{\linewidth}{3pt}\\%[1cm]  
{\bf {\vspace{-0.9cm} {\scriptsize Una Guía rápida a la programación con Python 3}}}


\vspace{2cm} 

\begin{center}
	{\ {\bf \vspace{0.2cm}  Trabajo presentado por el Profesor Linis B. Guerrero P.\\ 
			\vspace{0.2cm} 	
	}}
\end{center}




\begin{center}
	{\ {\bf \vspace{0.2cm}  		Autor: Linis B. Guerrero P.\\
			Traducido: Prof. Linis B. Guerrero P.\\
			\vspace{0.2cm} 	
	}}
\end{center}


\begin{center}
	{\ {\bf \vspace{0.2cm}  	MÉRIDA, VENEZUELA \\
			Septiembre, 2023 \\
			\vspace{0.2cm} 
	}}
\end{center}



\thispagestyle{empty}~~\vspace{-0.2cm}






%%%%%%%%%%%%%%%%%%%%%%%%%%%%%%%%%%%%%%%%%%%%%%%%%%%%%%%%%%%%%%%%%%%%%%%%%%%%%%%%%%%%%%%%%%%%%%%% 
%%                           Van Dos Pagina Universidad                                       %%
%%%%%%%%%%%%%%%%%%%%%%%%%%%%%%%%%%%%%%%%%%%%%%%%%%%%%%%%%%%%%%%%%%%%%%%%%%%%%%%%%%%%%%%%%%%%%%%%









% 	\hspace{2cm}

\vskip 5cm \rule{\linewidth}{1.3pt}\\%[0.5cm]
\begin{center}
	{\  \LARGE  {\bf \vspace{0.2cm}  Un curso intensivo de Python 3   \\				 
			\vspace{0.2cm} Una introducción rápida a la programación con Python \\ }}						
\end{center}		


\vskip 0.5cm \rule{\linewidth}{2.3pt}\\%[1cm]  
{\bf {\vspace{-0.9cm} {\scriptsize Una Guía rápida a la programación con Python 3}}}





\thispagestyle{empty}~~\vspace{-0.2cm}


%	\vspace{10cm}

\newpage





\pagenumbering{Roman}
\setcounter{page}{1}
\pagestyle{plain}





\chapter*{Agradecimientos} 


%   https://www.lawebdelprogramador.com/foros/TeX-Latex/964661-renombrar-Indice-de-Cuadros-y-Figuras.html#i1267809

\newpage
\tableofcontents
\addcontentsline{toc}{chapter}{Índice general}	
\newpage

\listoffigures
\addcontentsline{toc}{chapter}{Índice de figuras}	
\newpage

%	\tablename
\listoftables
\addcontentsline{toc}{chapter}{Índice de tablas Cuadros}	
\newpage










%	\newpage

\addtocontents{toc}{\hfill \textbf{} \par}

\addtocontents{toc}{\hspace{-7.5mm} \textbf{Capítulos}}

\addtocontents{toc}{\hfill \textbf{Página} \par}
\addtocontents{toc}{\vspace{-2mm} \hspace{-7.5mm} \hrule \par}
\addcontentsline{toc}{chapter}{Agradecimientos }	



%\thispagestyle{empty}~~\vspace{-0.2cm}

%\pagenumbering{Roman}
%\setcounter{page}{1}
%\pagestyle{plain}




%\thispagestyle{empty}~~\vspace{-0.2cm}


















%    instruciones comandos para tablas de contenidos  







\chapter*{Prefacio}	
\addcontentsline{toc}{chapter}{Prefacio}



En proceso .... 


El concepto de prefacio de un libro es una introducción escrita por el autor o por otra persona, que se coloca al principio del libro y tiene como objetivo proporcionar información sobre la obra y su contexto. 

El prefacio puede incluir detalles sobre la motivación del autor para escribir el libro, cómo surgió la idea, el proceso de creación, la investigación o la influencia de otras obras de arte o personas. También puede incluir una descripción del contenido del libro y cómo se relaciona con otros temas o trabajos similares. En resumen, el prefacio es una introducción al libro y su contexto que ayuda al lector a comprender mejor la obra.

Un prefacio es, en literatura, un texto de introducción y de presentación, ubicado al inicio de un libro. En el prefacio se da a conocer el plan y los puntos de vista utilizados durante la elaboración del escrito, a la par que allí también corresponde prevenir sobre posibles objeciones o reservas, responder a críticas ya formuladas a las ediciones anteriores o a los avances de la obra, y finalmente también dar ideas sobre el mensaje que el autor quiere transmitir con este documento; por ejemplo, y si el escrito planteara inquietudes sociales, el mensaje principal o más trascendente podría estar vinculado con la organización social, la pobreza, la desigual distribución de recursos, la prostitución, la violencia, el consumo de drogas y el narcotráfico, etc, o con varios de los temas citados.

El prefacio generalmente es corto cuando el mismo se orienta y se centra a ser una advertencia, y usualmente es largo cuando también incluye prolegómenos, motivaciones profundas o casuales, antecedentes, etc.


\section*{A Quién va dirigido este libro}
\addcontentsline{toc}{section}{A Quién va dirigido este libro}

El objetivo de este libro es enseñar de manera didáctica las herramientas fundamentales para programa en Python.   Está dirigido a personas de cualquier edad que deseen aprender a programar, y que tal vez nunca hayan programado en Python o nunca hayan programado en absoluto. Se hacen  comentarios  breve respecto al vocabulario común usado con los lenguajes de alto nivel en informática, para familiarizar al lector que se inicia en Python.  Dando las pautas para instalar Python y un entrono para editar  código. Se exponen los conceptos e instrucciones básicos de programación con ejemplos de código, para ir formando una base sólida en los conceptos  programación, habilidad y seguridad para su desarrollo como programador. Se trabajan ejemplos que sirvan de referencia para crear programas\footnote{Un programa es una serie de instrucciones ordenadas, codificadas en un lenguaje de programación que expresa uno o varios algoritmos y que puede ser ejecutado en una computadora.} como juegos, visualizaciones de datos, y aplicaciones web, mientras se desarrolla una base en programación que te servirá bien por el resto de tu vida.\\


El contenido del libro es ideal para profesores de todos los niveles que quieren ofrecer a sus estudiantes una introducción
a la programación basada en proyectos. Si estás tomando una clase universitaria y quieres una introducción  más amigable 
a Python que el texto que te han asignado, este libro puede hacer que tu clase sea más fácil. Si está buscando cambiar de carrera, 
este curso de Python puede ayudarlo hacer la transición hacia una carrera profesional más satisfactoria.



\section*{¿Qué esperas aprender en este libro?}
\addcontentsline{toc}{section}{¿Qué esperas aprender en este libro?}
El propósito de este libro es darle las herramientas para que se convierta en un buen programador de Python  y  en general en un buen  programador;  además,  aprenderás y adoptará buenos hábitos para desarrollar Python.

En la Parte I de este libro, aprenderá conceptos básicos de programación  necesarios para escribir programas en Python. Estos conceptos son casi los mismos que aprenderías al comenzar a estudiar cualquier lenguaje de programación.

Aquí también aprenderá sobre tipos de datos y las formas en que puede almacenarlos en sus programas. Crearás colecciones de 
datos, como listas y diccionarios, y trabajará en esas colecciones de manera eficiente. 
Aprenderás usar bucles {\ttfamily while} y declaraciones{\ttfamily if} para probar ciertas condiciones, por lo que puede ejecutar bloques específicos de código mientras esas condiciones sean verdaderas y se ejecuten otras en caso que no sean verdaderas, una técnica que le 
ayuda a automatizar muchas procesos.\\

Aprenderá a aceptar aportaciones de los usuarios para que sus programas sean interactivos y mantener sus programas funcionando todo 
el tiempo que el usuario lo requiera,  lo harás explorar cómo escribir funciones que hagan que partes de su programa sean reutilizables,
por lo que solo tienes que escribir bloques de código que realicen ciertas acciones una vez, mientras usas ese código tantas veces 
como necesites. Luego ampliarás este concepto a comportamientos más complicados con clases, haciendo que programas bastante simples respondan a una variedad de situaciones. Aprenderá a escribir programas que manejen correctamente los errores comunes.

Después de trabajar en cada uno de estos conceptos básicos, escribirá una serie de programas cada vez más complejos utilizando lo que ha aprendido. Finalmente, darás tu primer paso hacia la programación intermedia al aprender a escribir pruebas para tu código, de modo que puedas desarrollar más sus programas sin preocuparte por introducir errores. Toda la información de la Parte I lo preparará para asumir proyectos más grandes y complejos.\\


En la Parte II, en proceso... aplicará lo aprendido en la Parte I a tres proyectos. Puede realizar cualquiera o todos estos proyectos, en el orden que mejor le convenga. En los capítulos 12 a 14 está el primer proyecto,  un juego de disparos al estilo Space Invaders llamado Alien Invasion, que incluye varios niveles, cada vez más difíciles. Una vez que hayas completado este proyecto, debes estar en camino de poder desarrollar tus propios juegos en 2D. Incluso si no aspiras a convertirte en programador de juegos, trabajar en este proyecto es una forma divertida de poner a prueba  parte de lo que aprendió  en la Parte I.\\

El los capítulos del 15 al 17 se expone el segundo proyecto, en esta caso se trata de técnicas para visualización de datos. 
Los científicos de datos utilizan una variedad de técnicas de visualización para encontrar patrones y dar sentido a la gran cantidad de información disponible. Trabajará con conjuntos de datos conjuntos de datos que descargue de fuentes en línea y conjuntos de datos que sus programas descarguen automáticamente. Una vez que haya completado este proyecto, podrá escribir programas que analicen grandes conjuntos de datos y creen representaciones visuales de muchos tipos diferentes de información.\\

El tercer proyecto se presenta en los Capítulos del 18 al 20, creará una pequeña aplicación web llamada Registro de aprendizaje. Este proyecto le permite llevar un diario organizado de la información que ha aprendido sobre un tema específico. Podrá mantener registros separados para diferentes temas y permitir que otros creen una cuenta y comiencen sus propios diarios. También aprenderá cómo implementar su proyecto para que cualquiera pueda acceder a él en línea, desde cualquier parte del mundo.


\section*{Navegando en este libro}
\addcontentsline{toc}{section}{Navegando en este libro}

Este proyecto esta ideado en dos partes, la primera parte  enseña los conceptos básicos que necesitará para escribir programas en Python. Muchos de estos conceptos son comunes a todos los lenguajes de programación, por lo que te serán útiles a lo largo de tu vida como programador. Si ya tienes configurado Python y tienes un editor o IDE que te gusta, puedes comenzar en el capítulo \ref{Cap3} sin problemas.\\

Este libro está organizado de la siguiente manera:  

\begin{list}{$\bullet$}{}  
	
	\item En el \hyperlink{Cap1}{Capítulo \ref{Cap1}} se hace una introducción breve a Python, algunos los términos y conceptos comunes en informática, específicamente los que se asocian con Python, y una descripción de algunos editores e IDEs compatibles con Python.
		
	\item En el \hyperlink{Cap2}{Capítulo \ref{Cap2}} se inicia enseñando la instalación y configuración tanto para Python, como para algunos editores e IDEs, y se describen las características más importantes. Se ejecutará la primera su experiencia con un fragmento de código Python y su primer programa Python, con el interprete de Python tanto en la terminal de Windows (cmd), como en  el IDLE de Python o VSCode.
		
	
	\item En el \hyperlink{Cap3}{Capítulo \ref{Cap3}} se estudian los tipos de datos, y se asigna información a variables, es una buena lectura para empezar. Y se estudia lo referentes a cadenas o string.
	
	\item  En el \hyperlink{Cap4}{Capítulo \ref{Cap4}}  presentan las Tuplas, Listas, Diccionarios, conjuntos, y tipo  Colecciones. Las listas pueden almacenar tanta información como desee en un solo lugar, lo que le permitirá trabajar con esos datos de manera eficiente. Podrás trabajar con cientos, miles e incluso millones de valores en tan solo unas pocas líneas de código.
	En la \hyperlink{diccionario}{Sección \ref{diccionario}} se empieza a utilizar los diccionarios de Python, que le permiten establecer conexiones entre diferentes estructuras de información. Al igual que las listas, los diccionarios pueden contener tanta información como necesite almacenar.


	\item  En el \hyperlink{Cap5}{Capítulo \ref{Cap5}}  aprenderá cómo aceptar entradas de los usuarios para que sus programas sean interactivos. También aprenderá sobre los bucles{\ttfamily while}, que ejecutan bloques de código repetidamente siempre que ciertas condiciones sigan siendo verdaderas.
	
	
	\item En el \hyperlink{Cap6}{Capítulo \ref{Cap6}}  se introducen las estructuras de control de flujo: {\ttfamily if-elif-else}, para escribir código que responda de una manera si ciertas condiciones son verdaderas y responda de otra manera si esas condiciones no son verdaderas.
	
	% Cap7 = imput
	
	
	
	\item  En el \hyperlink{Cap7}{Capítulo \ref{Cap7}} escribirás funciones, que son bloques de código con nombre que realizan una tarea específica y se pueden ejecutar cuando las necesites.


   \item  En el \hyperlink{Cap8}{Capítulo \ref{Cap8}}  se inicia el trabajo con Módulos, Paquetes y programas,...

 
   \item  En el \hyperlink{Cap9}{Capítulo \ref{Cap9}} se presentan las clases, que le permiten modelar objetos del mundo real. Escribirás código que represente perros, gatos, personas, automóviles, cohetes y más.



  \item  En el \hyperlink{Cap10}{Capítulo \ref{Cap10}}  se inicia el trabajo archivos,  y se enseña a  manejar errores para que sus programas no colapsen inesperadamente. Almacenará datos antes de que se cierre su programa y los volverá a leer cuando el programa se ejecute nuevamente. 
Aprenderá sobre las excepciones de Python, que le permiten anticipar errores y hacer que sus programas los manejen correctamente.


	
	\item  En el \hyperlink{Cap11}{Capítulo \ref{Cap11}} aprenderá a escribir pruebas eficientes para comprobar que sus programas funcionan de la forma prevista. Como resultado, podrá ampliar sus programas sin preocuparse por introducir nuevos errores. Probar tu código es una de las primeras habilidades que te ayudarán a pasar de programador principiante a programador intermedio.
	
	
	
	\item  En proceso... El proyecto Invasión  alienígena en los \hyperlink{Cap12}{Capítulo \ref{Cap12}} y \ref{Cap13} incluye una configuración para controlar la velocidad  de fotogramas, lo que hace que el juego se ejecute de manera más consistente en diferentes sistemas operativos. Se utiliza un enfoque más simple para construir la flota de extraterrestres y también se ha limpiado la organización general del proyecto.

    \item   En proceso... Los proyectos de visualización de datos se desarrolla en los \hyperlink{Cap15}{Capítulo \ref{Cap17}} a \ref{Cap18}  utilizan las funciones más recientes de Matplotlib y Plotly. Las visualizaciones de Matplotlib presentan configuraciones de estilo actualizadas. El proyecto de caminata aleatoria tiene una pequeña mejora que aumenta la precisión de los gráficos, lo que significa que verá surgir una variedad más amplia de patrones cada vez que genere una nueva caminata. Todos los proyectos que presentan Plotly ahora utilizan el módulo Plotly Express, que le permite generar sus visualizaciones iniciales con solo unas pocas líneas de código. Puede explorar fácilmente una variedad de visualizaciones antes de comprometerse con un tipo de trama y luego concentrarse en refinar los elementos individuales de esa trama.

    \item  En proceso... El proyecto Learning Log se presenta en los \hyperlink{Cap11}{Capítulo \ref{Cap18}} a  \ref{Cap19} se creó utilizando la última versión de Django y se diseñó utilizando la última versión de Bootstrap. Se ha cambiado el nombre de algunas partes del proyecto para que sea más fácil seguir la organización general del proyecto. El proyecto ahora está implementado en Platform.sh, un moderno servicio de alojamiento para proyectos de Django. El proceso de implementación está controlado por archivos de configuración YAML, que le brindan un gran control sobre cómo se implementa su proyecto. Este enfoque es consistente con la forma en que los programadores profesionales implementan proyectos modernos de Django.

   \item  En proceso...El Apéndice xx dirige a los lectores a varios de los recursos en línea más populares para obtener ayuda. 
\end{list} 


Gracias por leer el curso intensivo de Python! Si tiene algún comentario o pregunta, no dude en ponerse en contacto;  Soy Julian en Celcie.Ula.ve  jajajaja   En proceso...

%Soy @ehmatthes en Twitter.




\section*{¿Por qué Python?}
\addcontentsline{toc}{section}{¿Por qué Python?}

Python es un lenguaje increíblemente eficiente: sus programas harán más en menos líneas de código de lo que requerirían muchos otros lenguajes. Es un lenguaje de programación interpretado cuya filosofía hace hincapié en una sintaxis muy limpia y un código legible, que le te ayudará a escribir código "limpio". Su código será más fácil de leer, más fácil de depurar y más fácil de ampliar y desarrollar, en comparación con otros lenguajes. 

Python es una muy buena alternativa para empezar a programar puesto que es un lenguaje muy sencillo, fácil y con una curva de aprendizaje buena. La sintaxis es fácil de entender puesto que es cercana al lenguaje natural.\\

La gente usa Python para muchos propósitos: crear juegos, crear aplicaciones web, resolver problemas comerciales y desarrollar herramientas internas en todo tipo de empresas interesantes, en el campo del aprendizaje automático y la inteligencia artificial. Actualmente se utiliza en campos científicos, para investigación académica y trabajos aplicados.\\

Una de las razones más importantes por las que recomiendo el uso de Python es por la comunidad Python, que incluye un grupo de personas increíblemente diversa y acogedora. La comunidad es esencial para los programadores porque la programación no es una actividad solitaria. La mayoría de nosotros, incluso los programadores más experimentados, necesitamos pedir consejo a otras personas que ya han resuelto problemas similares. Tener una comunidad bien conectada y de apoyo es fundamental para ayudarle a resolver problemas, y la comunidad de Python apoya plenamente a las personas que están aprendiendo Python como su primer lenguaje de programación o que llegan a Python con experiencia en otros lenguajes. Python es un gran lenguaje para aprender, ¡así que comencemos!






\section*{Recursos en línea}
\addcontentsline{toc}{section}{Recursos en línea}

Al estudiar este libro, puede visitar el sitio web   para obtener documentación detallada sobre  clases, funciones, y muchos ejemplos. También hay un curso con  vídeos creado por 


En proceso.....\\

No Starch Press tiene más información sobre este libro disponible en línea en 
% https://nostarch.com/python-crash-course-3rd-edition.
También mantengo un amplio conjunto de recursos complementarios en % https://ehmatthes.github.io/pcc_3e. 

Estos recursos incluyen lo siguiente:

Instrucciones de configuración Las instrucciones de configuración en línea son idénticas a las del libro, pero incluyen enlaces activos en los que puede hacer clic para seguir los diferentes pasos. Si tiene algún problema de configuración, consulte este recurso.
Actualizaciones Python, como todos los lenguajes, está en constante evolución. Mantengo un conjunto completo de actualizaciones, por lo que si algo no funciona, consulte aquí para ver si las instrucciones han cambiado.

Soluciones a los ejercicios Debe dedicar mucho tiempo a realizar los ejercicios de las secciones “Pruébelo usted mismo”. Sin embargo, si estás estancado y no puedes progresar, las soluciones para la mayoría de los ejercicios están en línea.
Hojas de referencia También está disponible en línea un conjunto completo de hojas de referencia descargables para una referencia rápida a los conceptos principales.

\section*{Ejemplos de código}
\addcontentsline{toc}{section}{Ejemplos de código}

En proceso...\\

El  material suplementario (ejemplos de código, cuadernos IPython, etc.) está disponible para su descarga en:


\hspace{0.5cm} \href{https://github.com/amueller/introduction_to_ml_with_python}{\color{vino} https://github.com/amueller/introduction\_to\_ml\_with\_python}.\\

Este libro se desarrolló con la finalidad de  ayudarle en  su trabajo. En general, el  código en los  ejemplo del libro lo puede usar en sus programas y documentación. No necesita contactar con nosotros para obtener permiso a menos que esté reproduciendo una parte significativa del código. Por ejemplo, escribir un programa que usa varios trozos de código de este libro no requiere permiso. Vender o distribuir un \\



Si cree que el uso que hace de los ejemplos de código queda fuera del uso legítimo o del permiso otorgado
arriba, no dude en contactarnos en {\color{vino} permissions@oreilly.com}. 


\section*{Convenciones utilizadas en este libro}
\addcontentsline{toc}{section}{Convenciones utilizadas en este libro}

En proceso..\\
En este libro se utilizan las siguientes convenciones tipográficas:\\

\noindent{\bf{Italic}}

\noindent Indica nuevos términos, direcciones URL, direcciones de correo electrónico, nombres de archivo y extensiones de archivo. \\

\noindent {\bf{Constant width}}

\noindent Usado para listados de programas, así como dentro de párrafos para referirse a elementos de programas tales como nombres de variables o funciones, bases de datos, tipos de datos, variables de entorno, sentencias y palabras clave. También se utiliza para comandos y nombres de módulos y paquetes. \\


\noindent {\bf Constant width bold}

\noindent Muestra comandos u otro texto que el usuario debe escribir literalmente.\\ 


\noindent {\bf Constant width italic}

\noindent Muestra texto que debe reemplazarse por valores proporcionados por el usuario o por valores determinados por el contexto. \\


\vspace{0.4cm}
\noindent \begin{minipage}[t]{0.2\textwidth}
	\centering\raisebox{\dimexpr \topskip-\height}{% 		
\includegraphics[scale=0.35]{../Imagen_Machine/Zorrillo}}	
	
\includegraphics[scale=0.25]{../Logos_de_Python/Sugerencia}}	
\end{minipage}%\hfill
\begin{minipage}[t]{0.75\textwidth} 
	
	Este elemento significa sugerencia.
\end{minipage}
\vspace{0.3cm}



\vspace{0.2cm}
\noindent \begin{minipage}[t]{0.2\textwidth}
	\centering\raisebox{\dimexpr\topskip-\height}{% 
\includegraphics[scale=0.35]{../Imagen_Machine/Pajaro}}  %[width=\textwidth]
		
\includegraphics[scale=0.20]{../Logos_de_Python/Busquedas_R}}
\end{minipage}%\hfill
\begin{minipage}[t]{0.75\textwidth} 
	
	Este elemento significa una nota general.
\end{minipage}
\vspace{0.2cm}



\vspace{0.3cm}
\noindent\begin{minipage}[t]{0.2\textwidth}
	\centering\raisebox{\dimexpr \topskip-\height}{%		
\includegraphics[scale=0.35]{../Imagen_Machine/Escorpion}}  
	\includegraphics[scale=0.20]{../Logos_de_Python/Advertencia1}} 
\end{minipage}%\hfill
\begin{minipage}[t]{0.70\textwidth} 
	
	Este icono indica una advertencia o precaución.
\end{minipage} 
\vspace{0.5cm}




\section*{Libros online Ula}
\addcontentsline{toc}{section}{Libros online Ula}

En proceso..\\



\section*{Como contactar con nosotros}
\addcontentsline{toc}{section}{Como contactar con nosotros}

En proceso...\\

Dirija sus comentarios y preguntas sobre el libro al editor:\\


O’Reilly Media, Inc.

1005 Gravenstein Highway North

Sebastopol, CA 95472

800-998-9938 (in the United States or Canada)

707-829-0515 (international or local)

707-829-0104 (fax)\\




Tenemos una página web para este libro, con la fe de erratas, ejemplos y cualquier información adicional. Puede acceder a esta página en:\\ \href{http://bit.ly/intro-machine-learning-python}{\color{vino}http://bit.ly/intro-machine-learning-python}.\\

Para comentar o hacer preguntas técnicas sobre este libro, envíe un correo electrónico a {\color{vino} bookquestions@oreilly.com}.\\

Para obtener más información sobre nuestros libros, cursos, conferencias y noticias, consulte nuestro sitio web en \href{https://www.oreilly.com/}{\color{vino} http://www.oreilly.com}.\\

%\dfrac{núm}{den}Encuéntrenos en Facebook:   \href{ http://facebook.com/oreilly}{\color{vino} http://facebook.com/oreilly}.\\
Síganos en Twitter:   \href{ http://twitter.com/oreillymedia}{\color{vino} http://twitter.com/oreillymedia}.\\
Míranos en YouTube:   \hspace{0.5cm} \href{http://www.youtube.com/oreillymedia}{\color{vino} http://www.youtube.com/oreillymedia}.\\


\section*{Agradecimientos}
\addcontentsline{toc}{section}{Agradecimientos}


{\bf{De Linis}}


Sin la ayuda y el apoyo de un gran grupo de personas, este libro nunca habría existido.\\

Me gustaría agradecer a los editores, 

Estoy eternamente agradecido por la acogedora comunidad científica de código abierto de Python, 
especialmente los contribuyentes 

También estoy agradecido por las discusiones con muchos de mis colegas y compañeros que me ayudaron
a comprender los desafíos del aprendizaje automático y me dio ideas para la estructura un libro de texto. \\


En el aspecto personal, quiero agradecer a mis padres,  por su continuo apoyo y aliento. También quiero agradecer a muchas personas en mi 
vida cuyo amor y amistad me dieron la energía y el apoyo para emprender una tarea tan desafiante.\\



{\bf{De Julian}}


Me gustaría agradecer, sin cuya ayuda y orientación este proyecto
ni siquiera habría existido. Gracias a Carlos Dávila, Oswaldo  por leerlo en los primeros días. Gracias a la gente ULA Ciencias por su infinita paciencia. Y finalmente gracias a  Dios , por su apoyo eterno e infinito.







\addcontentsline{toc}{chapter}{Prefacio}


\setcounter{secnumdepth}{3}



\setcounter{tocdepth}{3}






%















